% ================================================================
%  preamble.tex — the book's design. Everything visual lives here.
%  Trim: 7.5 x 9.25 in (classic technical-book size).
% ================================================================

% ---- Page geometry ----------------------------------------------
\usepackage[paperwidth=7.5in, paperheight=9.25in,
            inner=0.85in, outer=0.9in, top=0.95in, bottom=1in,
            headheight=14pt, headsep=0.28in, footskip=0.45in]{geometry}

% ---- Fonts --------------------------------------------------------
\usepackage{fontspec}
\defaultfontfeatures[\rmfamily,\sffamily]{Ligatures=TeX}
\setmainfont{LibertinusSerif}[
  Extension=.otf, UprightFont=*-Regular, BoldFont=*-Semibold,
  ItalicFont=*-Italic, BoldItalicFont=*-SemiboldItalic]
\IfFontExistsTF{FiraSans-Regular.otf}{%
  \setsansfont{FiraSans}[
    Extension=.otf, UprightFont=*-Regular, BoldFont=*-SemiBold,
    ItalicFont=*-Italic, BoldItalicFont=*-SemiBoldItalic, Scale=0.92]
}{%
  \setsansfont{LibertinusSans}[
    Extension=.otf, UprightFont=*-Regular, BoldFont=*-Bold, ItalicFont=*-Italic]
}
% Menlo: macOS system mono with full box-drawing coverage for diagrams.
\setmonofont{Menlo}[Scale=0.85]
\newfontfamily\symbolfont{DejaVuSans.ttf}[BoldFont=DejaVuSans-Bold.ttf, ItalicFont=DejaVuSans-Oblique.ttf]

% Glyphs missing from Libertinus -> DejaVu fallback.
\usepackage{newunicodechar}
% Inside code/diagrams (Menlo) keep the monospace glyph so columns stay aligned;
% only proportional text falls back to DejaVu Sans.
\makeatletter
\newcommand{\sym}[1]{\begingroup\edef\@tempa{\f@family}\edef\@tempb{\ttdefault}%
  \ifx\@tempa\@tempb {\upshape #1}\else{\symbolfont #1}\fi\endgroup}
\makeatother
\newunicodechar{✓}{\sym{✓}}
\newunicodechar{✔}{\sym{✔}}
\newunicodechar{✗}{\sym{✗}}
\newunicodechar{►}{\sym{►}}
\newunicodechar{◄}{\sym{◄}}
% Emoji (colour, rendered to PNG from Apple Color Emoji) and CJK
\newcommand{\emoji}[1]{\raisebox{-0.12em}{\includegraphics[height=0.95em]{assets/emoji/#1.png}}}
\newunicodechar{🍵}{\emoji{1f375}}
\newunicodechar{🍰}{\emoji{1f370}}
\newunicodechar{🎉}{\emoji{1f389}}
\newunicodechar{✨}{\emoji{2728}}
\newunicodechar{❌}{\emoji{274c}}
\newunicodechar{✅}{\emoji{2705}}
\newfontfamily\cjkfont{Hiragino Sans GB}[Scale=0.85]
\newunicodechar{你}{{\cjkfont 你}}
\newunicodechar{好}{{\cjkfont 好}}

% ---- Colour palette -----------------------------------------------
\usepackage[table]{xcolor}
\definecolor{brand}{HTML}{1F3A5F}      % deep navy
\definecolor{brandlight}{HTML}{E8EEF5}
\definecolor{accent}{HTML}{C8553D}     % terracotta
\definecolor{accent2}{HTML}{2A7F8E}    % teal
\definecolor{codebg}{HTML}{F6F8FA}
\definecolor{codeframe}{HTML}{D5DCE3}
\definecolor{muted}{HTML}{6B7280}

% ---- Core packages --------------------------------------------------
\usepackage{amssymb}        % task-list boxes (pandoc: $\square$, $\boxtimes$)
\usepackage{microtype}
\usepackage{graphicx}
\usepackage{calc}
\usepackage{array}
\usepackage{longtable}
\usepackage{booktabs}
\usepackage{etoolbox}
\usepackage{enumitem}
\usepackage{fvextra}
\usepackage{tikz}
\usetikzlibrary{positioning,calc,arrows.meta,shapes.geometric}
\usepackage[most]{tcolorbox}
\usepackage{titlesec}
\usepackage{titletoc}
\usepackage{fancyhdr}
\usepackage{emptypage}
\usepackage{needspace}

\raggedbottom
\frenchspacing
\setlength{\parindent}{1.2em}
\setlength{\parskip}{0pt}
\setlength{\emergencystretch}{3em}
\tolerance=1200
\linespread{1.06}
\setlist{itemsep=2pt, topsep=5pt, parsep=0pt, leftmargin=1.6em}
\setlist[itemize]{label=\textcolor{accent}{\small\textbullet}}
\providecommand{\tightlist}{\setlength{\itemsep}{1pt}\setlength{\parskip}{0pt}}

% Tables (pandoc emits longtable + booktabs)
\newcounter{none}   % pandoc sets \LTcaptype{none} for uncaptioned tables (newer longtable needs it)
\renewcommand{\arraystretch}{1.3}
\BeforeBeginEnvironment{longtable}{\Needspace{9\baselineskip}}   % \Needspace: plain \needspace can print the header twice
\AtBeginEnvironment{longtable}{\small}
\setlength{\LTpre}{10pt}
\setlength{\LTpost}{10pt}
\arrayrulecolor{brand!70}

% ---- Code: pandoc highlighting + our boxes --------------------------
\DefineVerbatimEnvironment{Highlighting}{Verbatim}{commandchars=\\\{\}}
\newcommand{\AlertTok}[1]{\textcolor[rgb]{0.94,0.16,0.16}{#1}}
\newcommand{\AnnotationTok}[1]{\textcolor[rgb]{0.56,0.35,0.01}{\textbf{\textit{#1}}}}
\newcommand{\AttributeTok}[1]{\textcolor[rgb]{0.13,0.29,0.53}{#1}}
\newcommand{\BaseNTok}[1]{\textcolor[rgb]{0.00,0.00,0.81}{#1}}
\newcommand{\BuiltInTok}[1]{#1}
\newcommand{\CharTok}[1]{\textcolor[rgb]{0.31,0.60,0.02}{#1}}
\newcommand{\CommentTok}[1]{\textcolor[rgb]{0.56,0.35,0.01}{\textit{#1}}}
\newcommand{\CommentVarTok}[1]{\textcolor[rgb]{0.56,0.35,0.01}{\textbf{\textit{#1}}}}
\newcommand{\ConstantTok}[1]{\textcolor[rgb]{0.56,0.35,0.01}{#1}}
\newcommand{\ControlFlowTok}[1]{\textcolor[rgb]{0.13,0.29,0.53}{\textbf{#1}}}
\newcommand{\DataTypeTok}[1]{\textcolor[rgb]{0.13,0.29,0.53}{#1}}
\newcommand{\DecValTok}[1]{\textcolor[rgb]{0.00,0.00,0.81}{#1}}
\newcommand{\DocumentationTok}[1]{\textcolor[rgb]{0.56,0.35,0.01}{\textbf{\textit{#1}}}}
\newcommand{\ErrorTok}[1]{\textcolor[rgb]{0.64,0.00,0.00}{\textbf{#1}}}
\newcommand{\ExtensionTok}[1]{#1}
\newcommand{\FloatTok}[1]{\textcolor[rgb]{0.00,0.00,0.81}{#1}}
\newcommand{\FunctionTok}[1]{\textcolor[rgb]{0.13,0.29,0.53}{\textbf{#1}}}
\newcommand{\ImportTok}[1]{#1}
\newcommand{\InformationTok}[1]{\textcolor[rgb]{0.56,0.35,0.01}{\textbf{\textit{#1}}}}
\newcommand{\KeywordTok}[1]{\textcolor[rgb]{0.13,0.29,0.53}{\textbf{#1}}}
\newcommand{\NormalTok}[1]{#1}
\newcommand{\OperatorTok}[1]{\textcolor[rgb]{0.81,0.36,0.00}{\textbf{#1}}}
\newcommand{\OtherTok}[1]{\textcolor[rgb]{0.56,0.35,0.01}{#1}}
\newcommand{\PreprocessorTok}[1]{\textcolor[rgb]{0.56,0.35,0.01}{\textit{#1}}}
\newcommand{\RegionMarkerTok}[1]{#1}
\newcommand{\SpecialCharTok}[1]{\textcolor[rgb]{0.81,0.36,0.00}{\textbf{#1}}}
\newcommand{\SpecialStringTok}[1]{\textcolor[rgb]{0.31,0.60,0.02}{#1}}
\newcommand{\StringTok}[1]{\textcolor[rgb]{0.31,0.60,0.02}{#1}}
\newcommand{\VariableTok}[1]{\textcolor[rgb]{0.00,0.00,0.00}{#1}}
\newcommand{\VerbatimStringTok}[1]{\textcolor[rgb]{0.31,0.60,0.02}{#1}}
\newcommand{\WarningTok}[1]{\textcolor[rgb]{0.56,0.35,0.01}{\textbf{\textit{#1}}}}

% Kate marks WARNING/TODO/TESTING in comments as "alerts" (red) — show them as comments.
\renewcommand{\AlertTok}[1]{\CommentTok{#1}}
\newcommand{\codesize}{\fontsize{9}{11.2}\selectfont}
\newenvironment{Shaded}{}{}
\RecustomVerbatimEnvironment{Highlighting}{Verbatim}{%
  commandchars=\\\{\}, fontsize=\codesize, breaklines, breakanywhere, breaknonspaceingroup,
  breakanywheresymbolpre={}, breaksymbolleft={\tiny\textcolor{muted}{$\hookrightarrow$}}}
\RecustomVerbatimEnvironment{verbatim}{Verbatim}{%
  fontsize=\codesize, breaklines, breakanywhere, breakanywheresymbolpre={},
  breaksymbolleft={\tiny\textcolor{muted}{$\hookrightarrow$}}}

\tcbset{bookbox/.style={enhanced, breakable, vfill before first=false, arc=1.5pt, boxrule=0.5pt,
  before skip=9pt, after skip=11pt},
  short/.style={unbreakable}}   % filter.lua passes "short" for boxes that should never split

% Source code (has a language); codeboxcap adds a "Listing X.Y" caption bar.
\tcbset{codestyle/.style={bookbox, colback=codebg, colframe=codeframe,
  borderline west={2.5pt}{0pt}{accent2},
  left=8pt, right=6pt, top=3pt, bottom=3pt}}
\newtcolorbox{codebox}[1][]{codestyle, #1}
\newtcolorbox{codeboxtight}[1][]{codestyle, before skip=3pt, #1}
% Keep a caption with its listing. \needspace is not used here: it inserts \penalty-100,
% which lets TeX break the page right after a heading and strand it.
\makeatletter
\newcommand{\keepspace}[1]{\par
  \ifdim\dimexpr\pagegoal-\pagetotal\relax<#1\relax
    \ifdim\pagegoal<\maxdimen\newpage\fi\fi}
\makeatother
\newenvironment{codeboxcap}[2][]{%
  \par\addvspace{10pt}\keepspace{6\baselineskip}\noindent
  {\sffamily\footnotesize\bfseries\color{brand}#2}\par\nobreak
  \begin{codeboxtight}[#1]}{\end{codeboxtight}}

% Terminal output / plain text (no language).
\newtcolorbox{outputbox}[1][]{bookbox, colback=white, colframe=codeframe,
  left=8pt, right=6pt, top=3pt, bottom=3pt, #1}

% ASCII / box-drawing diagrams — never line-broken, but may split across pages.
% #1 = box width (filter.lua computes it from the longest line) -> box is centred.
\newtcolorbox{diagrambox}[2][]{enhanced, breakable, arc=1.5pt, boxrule=0.4pt,
  colback=brandlight!35, colframe=codeframe,
  width=#2, enlarge left by={(\linewidth-#2)/2},
  left=6pt, right=6pt, top=6pt, bottom=6pt, before skip=10pt, after skip=12pt,
  #1}

% Callouts
\newtcolorbox{notebox}{bookbox, colback=accent2!6, colframe=accent2!6,
  borderline west={3pt}{0pt}{accent2}, boxrule=0pt, arc=0pt,
  left=10pt, right=8pt, top=6pt, bottom=6pt, fontupper=\small, unbreakable}
% Warnings ("**Warning:** …" paragraphs)
\newtcolorbox{warnbox}{bookbox, colback=accent!7, colframe=accent!7,
  borderline west={3pt}{0pt}{accent}, boxrule=0pt, arc=0pt,
  left=10pt, right=8pt, top=6pt, bottom=6pt, fontupper=\small, unbreakable}

\newtcolorbox{objectives}{bookbox, colback=brandlight!60, colframe=brand,
  boxrule=0.6pt, title={In this chapter},
  fonttitle=\sffamily\small\bfseries, colbacktitle=brand, coltitle=white,
  left=10pt, right=10pt, top=6pt, bottom=6pt, fontupper=\small\raggedright, unbreakable}

\newtcolorbox{takeaways}{bookbox, colback=accent!5, colframe=accent,
  boxrule=0.6pt, title={Key Takeaways},
  fonttitle=\sffamily\small\bfseries, colbacktitle=accent, coltitle=white,
  left=10pt, right=10pt, top=6pt, bottom=6pt, fontupper=\small\raggedright}

% ---- Headings -------------------------------------------------------
\setcounter{secnumdepth}{1}   % chapters + sections numbered
\setcounter{tocdepth}{1}

\titleformat{\chapter}[display]
  {\normalfont\sffamily\color{brand}}
  {\footnotesize\bfseries\color{accent}\addfontfeature{LetterSpace=18}CHAPTER\ \thechapter}
  {6pt}
  {\fontsize{26}{31}\selectfont\bfseries\raggedright\hyphenpenalty=10000 }
  [\vspace{2pt}]
\titleformat{name=\chapter,numberless}[display]
  {\normalfont\sffamily\color{brand}}{}{0pt}
  {\fontsize{26}{31}\selectfont\bfseries\raggedright\hyphenpenalty=10000 }
  [\vspace{2pt}]
\titlespacing*{\chapter}{0pt}{24pt}{18pt}

\titleformat{\section}
  {\normalfont\sffamily\Large\bfseries\color{brand}}
  {\textcolor{accent}{\thesection}}{0.7em}{}
\titleformat{\subsection}
  {\normalfont\sffamily\large\bfseries\color{brand}}{}{0em}{}
\titleformat{\subsubsection}
  {\normalfont\sffamily\normalsize\bfseries\color{brand!85}}{}{0em}{}
\titlespacing*{\section}{0pt}{18pt plus 3pt minus 2pt}{7pt}
\titlespacing*{\subsection}{0pt}{13pt plus 2pt minus 2pt}{5pt}
\titlespacing*{\subsubsection}{0pt}{10pt plus 2pt}{3pt}

% Subtitle line + rule under a chapter title
\newcommand{\chaptersubtitle}[1]{%
  \vspace*{-10pt}{\large\itshape\color{muted}#1\par}\vspace{6pt}}
\newcommand{\chapterrule}{%
  \noindent{\color{brand}\rule{\linewidth}{1.1pt}}\par\vspace{14pt}}

% Unnumbered structure used by Parts I–III and by un-numbered sections
\newcommand{\unnumberedchapter}[1]{%
  \cleardoublepage\phantomsection
  \addcontentsline{toc}{chapter}{#1}%
  \chapter*{#1}\markboth{#1}{#1}}
\newcommand{\unnumberedsection}[2]{% #1 = add to toc (true/false), #2 = title
  \section*{#2}\phantomsection
  \ifstrequal{#1}{true}{\addcontentsline{toc}{section}{#2}}{}%
  \markright{#2}}

% ---- Running heads -------------------------------------------------
\pagestyle{fancy}
\fancyhf{}
\newcommand{\headfont}{\sffamily\footnotesize\color{muted}}
\fancyhead[LE]{\headfont\thepage\hspace{1.2em}\leftmark}
\fancyhead[RO]{\headfont\rightmark\hspace{1.2em}\thepage}
\renewcommand{\headrulewidth}{0pt}
\renewcommand{\chaptermark}[1]{\markboth{Chapter \thechapter\enspace·\enspace #1}{}}
\renewcommand{\sectionmark}[1]{\markright{\thesection\enspace #1}}
\fancypagestyle{plain}{\fancyhf{}\fancyfoot[C]{\headfont\thepage}}

% ---- Table of contents styling -------------------------------------
\titlecontents{part}[0pt]
  {\addvspace{16pt}\sffamily\bfseries\color{brand}\large}
  {}{}{\hfill\contentspage}[\addvspace{3pt}]
\titlecontents{chapter}[2.3em]
  {\addvspace{5pt}\sffamily\bfseries}
  {\contentslabel{2.3em}}{\hspace*{-2.3em}}
  {\hfill\contentspage}
\titlecontents{section}[4.8em]
  {\small}
  {\contentslabel{2.5em}}{\hspace*{-2.5em}}
  {\titlerule*[0.7pc]{.}\contentspage}
% Room for four-digit page numbers (the book passes 1,000 pages)
\makeatletter
\renewcommand{\@pnumwidth}{2.2em}
\renewcommand{\@tocrmarg}{3.2em}
\makeatother
\contentsmargin{2.2em}
\newcommand{\booktocline}[1]{%
  \par\addvspace{10pt}\needspace{4\baselineskip}\noindent
  {\sffamily\scriptsize\bfseries\color{accent}\addfontfeature{LetterSpace=10}#1}\par
  \nobreak\addvspace{2pt}}

% ---- Part pages (the four lifecycle phases) ------------------------
% \lifecyclepart{<title>}{<description>}{<phase 1-4>}
\newcommand{\phasestrip}[1]{%
  \foreach \n/\lab in {1/Discover, 2/Plan, 3/Code, 4/Production}{%
    \ifnum\n=#1\relax
      \node[circle, fill=accent, text=white, minimum size=24pt,
            font=\sffamily\bfseries] (p\n) at (\n*2.4, 0) {\n};
      \node[below=4pt of p\n, font=\sffamily\scriptsize\bfseries, text=white] (q\n) {\lab};
      \node[above=5pt of p\n, font=\sffamily\tiny\bfseries\addfontfeature{LetterSpace=15},
            text=accent] {YOU ARE HERE};
    \else
      \node[circle, draw=white!55, text=white!70, minimum size=24pt,
            line width=0.8pt, font=\sffamily] (p\n) at (\n*2.4, 0) {\n};
      \node[below=4pt of p\n, font=\sffamily\scriptsize, text=white!60] (q\n) {\lab};
    \fi
  }
  \foreach \a/\b in {1/2, 2/3, 3/4}{%
    \draw[-{Latex[length=5pt]}, white!55, line width=0.8pt] (p\a) -- (p\b);}
  % the loop: what production teaches goes back to discovery
  \draw[-{Latex[length=5pt]}, white!45, line width=0.7pt]
    (q4.south) .. controls ([yshift=-0.7cm]q4.south) and ([yshift=-0.7cm]q1.south) .. (q1.south)
    node[pos=0.5, below=1pt, font=\sffamily\scriptsize\itshape, text=white!60] {learn};
}
% \prebookpart{<title>}{<description>} -- a front section, NOT one of the four parts.
% It does not advance the part counter, so Product Discovery stays Part I.
\newcommand{\prebookpart}[2]{%
  \cleardoublepage
  \phantomsection
  \addcontentsline{toc}{part}{#1}%
  \markboth{}{}%
  \thispagestyle{empty}%
  \begin{tikzpicture}[remember picture, overlay]
    \fill[brand] (current page.south west) rectangle (current page.north east);
    \node[anchor=north west, text=white!70, font=\sffamily\small\bfseries\addfontfeature{LetterSpace=25}]
      at ([xshift=0.9in, yshift=-3.0in]current page.north west) {BEFORE PART I};
    \node[anchor=north west, text=white, text width=5.4in, align=left,
          font=\sffamily\bfseries\fontsize{30}{36}\selectfont]
      at ([xshift=0.85in, yshift=-3.3in]current page.north west) {#1};
    \node[anchor=north west, text=white!80, text width=5in, align=left,
          font=\large\itshape]
      at ([xshift=0.9in, yshift=-4.75in]current page.north west) {#2};
    \begin{scope}[shift={([xshift=0.35in, yshift=1.35in]current page.south west)}]
      \phasestrip{0}
      \node[anchor=west, font=\sffamily\scriptsize, text=white!60] at (1.95, -2.6)
        {The whole map: \emph{How Software Gets Built}, page~\pageref*{ch:how-software-gets-built}};
    \end{scope}
  \end{tikzpicture}%
  \cleardoublepage}

\newcommand{\lifecyclepart}[3]{%
  \cleardoublepage
  \refstepcounter{part}%
  \phantomsection
  \addcontentsline{toc}{part}{Part \thepart\enspace\textperiodcentered\enspace #1}%
  \markboth{}{}%
  \thispagestyle{empty}%
  \begin{tikzpicture}[remember picture, overlay]
    \fill[brand] (current page.south west) rectangle (current page.north east);
    \node[anchor=north west, text=accent, font=\sffamily\bfseries\fontsize{110}{110}\selectfont]
      at ([xshift=0.85in, yshift=-1.1in]current page.north west) {\thepart};
    \node[anchor=north west, text=white!70, font=\sffamily\small\bfseries\addfontfeature{LetterSpace=25}]
      at ([xshift=0.9in, yshift=-3.0in]current page.north west) {PART \thepart};
    \node[anchor=north west, text=white, text width=5.4in, align=left,
          font=\sffamily\bfseries\fontsize{30}{36}\selectfont]
      at ([xshift=0.85in, yshift=-3.3in]current page.north west) {#1};
    \node[anchor=north west, text=white!80, text width=5in, align=left,
          font=\large\itshape]
      at ([xshift=0.9in, yshift=-4.75in]current page.north west) {#2};
    \begin{scope}[shift={([xshift=0.35in, yshift=1.35in]current page.south west)}]
      \phasestrip{#3}
      \node[anchor=west, font=\sffamily\scriptsize, text=white!60] at (1.95, -2.6)
        {The whole map: \emph{How Software Gets Built}, page~\pageref*{ch:how-software-gets-built}};
    \end{scope}
  \end{tikzpicture}%
  \cleardoublepage}

% ---- The six-step map, used on the stage pages ----------------------
% \mapstrip{<highlighted steps, e.g. 4 or 3,4>}   1 Discover … 6 Grow
\newcommand{\mapstrip}[1]{%
  \begin{tikzpicture}[x=0.8in]
    \foreach \n/\lab in {1/Discover, 2/Plan, 3/Code, 4/Ship, 5/Run, 6/Grow}{%
      \gdef\maphl{0}%
      \foreach \h in {#1}{\ifnum\h=\n\relax\gdef\maphl{1}\fi}%
      \ifnum\maphl=1\relax
        \node[circle, fill=accent, text=white, minimum size=16pt, inner sep=0pt,
              font=\sffamily\scriptsize\bfseries] (m\n) at (\n, 0) {\n};
        \node[below=3pt of m\n, font=\sffamily\scriptsize\bfseries, text=accent] (l\n) {\lab};
      \else
        \node[circle, draw=brand!35, text=brand!55, minimum size=16pt, inner sep=0pt,
              line width=0.7pt, font=\sffamily\scriptsize] (m\n) at (\n, 0) {\n};
        \node[below=3pt of m\n, font=\sffamily\scriptsize, text=muted] (l\n) {\lab};
      \fi
    }
    \foreach \a/\b in {1/2, 2/3, 3/4, 4/5, 5/6}{%
      \draw[-{Latex[length=4pt]}, brand!35, line width=0.6pt] (m\a) -- (m\b);}
    % steps 4–6 are the production phase
    \draw[brand!35, line width=0.6pt]
      ([yshift=9pt]m4.north west) -- ++(0, 3pt) -| ([yshift=9pt]m6.north east);
    \node[font=\sffamily\tiny\bfseries\addfontfeature{LetterSpace=15}, text=brand!55]
      at ([yshift=19pt]m5.north) {PRODUCTION};
    % the loop back
    \draw[-{Latex[length=4pt]}, brand!35, line width=0.6pt]
      (l6.south) .. controls ([yshift=-0.55cm]l6.south) and ([yshift=-0.55cm]l1.south) .. (l1.south)
      node[pos=0.5, below=0pt, font=\sffamily\scriptsize\itshape, text=muted] {learn};
  \end{tikzpicture}}

% ---- Stage pages (the eleven stages inside Part IV) -------------
% \bookpart{<n>}{<title>}{<the question it answers>}{<chapters>}{<map steps>}{<the piece>}
\newcommand{\bookpart}[6]{%
  \cleardoublepage
  \phantomsection
  \addtocontents{toc}{\protect\booktocline{STAGE #1\enspace·\enspace #2}}%
  \markboth{}{}%
  \thispagestyle{empty}%
  \begin{tikzpicture}[remember picture, overlay]
    \fill[brandlight] ([yshift=-3.6in]current page.north west) rectangle (current page.north east);
    \fill[accent] ([yshift=-3.6in]current page.north west) rectangle ([yshift=-3.66in]current page.north east);
    \node[anchor=south west, text=brand!25, font=\sffamily\bfseries\fontsize{150}{150}\selectfont]
      at ([xshift=0.7in, yshift=-3.55in]current page.north west) {#1};
  \end{tikzpicture}%
  \vspace*{2.9in}
  {\sffamily\small\bfseries\color{accent}\addfontfeature{LetterSpace=20}STAGE #1\par}
  \vspace{6pt}
  {\sffamily\bfseries\fontsize{26}{31}\selectfont\color{brand}#2\par}
  \vspace{16pt}
  {\color{brand}\rule{1.2in}{1pt}\par}
  \vspace{12pt}
  {\large\itshape\color{muted}#3\par}
  \vfill
  {\sffamily\scriptsize\bfseries\color{accent}\addfontfeature{LetterSpace=20}YOU ARE HERE\par}
  \vspace{8pt}
  \mapstrip{#5}\par
  \vspace{6pt}
  {\sffamily\small\color{brand}\textbf{The piece:} #6\par}
  \vspace{3pt}
  {\sffamily\footnotesize\color{muted}The whole map: \emph{How Software Gets Built}, page~\pageref*{ch:how-software-gets-built}\par}
  \vspace{26pt}
  {\sffamily\footnotesize\color{muted}#4\par}
  \cleardoublepage}

% ---- Links & PDF metadata -----------------------------------------
\usepackage{hyperref}
\usepackage{bookmark}
\usepackage{xurl}          % long URLs may break anywhere
\hypersetup{
  colorlinks=true, linkcolor=brand, urlcolor=accent2, citecolor=brand,
  pdftitle={The Complete Software Engineering Guide: From Deciding What to Build to a Running Production System},
  pdfsubject={Product discovery, project planning, writing good code, and building real applications: a Flask Notes API taken from a laptop to a monitored production service. First Edition.},
  pdfkeywords={software engineering, backend, Python, Flask, PostgreSQL, Docker, Nginx, CI/CD, deployment, observability, reliability, security, production systems},
  pdfcreator={The Complete Software Engineering Guide, First Edition},
  pdflang={en},
  pdfdisplaydoctitle=true,   % viewers show the title, not the file name
  bookmarksnumbered=true, pdfpagemode=UseOutlines}
